% !TEX root = ../main.tex

\chapter{试验台与硬件系统}
\chaptereng{Experimental Setup and Hardware System}


\section{新型非预混平面绝热火焰炉}

在项目初期，炉体制造使用“不锈钢毛细管插管钎焊”方案。该方案的具体工艺为：在无氧铜底座上钻出盲孔阵列，然后将外径与盲孔直径匹配的不锈钢毛细管插入盲孔底部，依靠钎焊密封固定。燃料（H$_2$）与氧化剂（O$_2$）分别经由不同的毛细管独立输送，在管口出口处相遇并形成非预混扩散火焰。

然而，在实际加工与前期点火测试中，该方案暴露出三重结构性缺陷：

\begin{enumerate}
	\item \textbf{钎焊工艺瓶颈}：为实现均匀温度场，需要尽可能缩小燃料管与氧化剂管之间的中心距（即管间距）。当管间距缩小至1.6mm以下时，相邻毛细管外壁之间的间隙就会非常小，钎焊时焊料难以均匀浸润，极易造成相邻管路之间的焊料搭接甚至堵塞，导致成品率急剧下降。
	\item \textbf{无主动冷却通道}：该加工方式需要在底座上先钻盲孔再插入毛细管，炉体内部空间已被气流通道完全占据，无法在不破坏气流分配均匀性的前提下额外加工水冷孔道。当燃烧甲烷/空气混合物（绝热温度$\sim$2200K）时，依靠自然对流和辐射散热尚可维持短时间（$<$30分钟）运行；但当燃烧H$_2$/O$_2$（绝热温度$>$3000K）时，顶部微管出口区域的表面热流密度高达30 MW/m$^2$以上，炉体顶部温度迅速超过铜的再结晶温度（$\sim$400$^\circ$C），导致材料强度下降、微管变形甚至烧穿。
\end{enumerate}

上述问题的根本原因在于“插焊”这一工艺路线本身的局限性。因此，需要彻底摒弃插管方案，采用新的加工思路重新设计炉体结构。

\subsection{新结构设计}

基于上述问题，本文设计的新结构可以直接由整块材料基体，直接利用数控深孔钻加工出微孔阵列及水冷通道，无需任何焊接或插接工序，避免了插管焊接导致的无法进一步缩小管间距和添加水冷通道的问题。

该加工方式带来的根本性改进体现在两个方面：

\begin{enumerate}
	\item \textbf{排布密度的提升}：由于取消了毛细管的外壁厚度和钎焊间隙的约束，孔与孔之间的最小中心距仅由钻孔工艺的最小壁厚决定。在采用电火花孔钻（孔径0.4mm，长径比$>$30:1）的条件下，最小壁厚可控制在0.4mm左右，对应管间距可缩小至1.6mm，为实现更高空间分辨率的均匀温度场创造了条件。
	\item \textbf{水冷通道的集成}：在微孔阵列外围及相邻气流通道之间的区域可以添加水冷通道，在不破坏气流通道的前提下，沿轴向钻设贯通水冷孔道（直径0.6mm）。水冷孔道分布于火焰面下方3--5mm处，冷却水强制对流可对炉体顶部壁面温度进行控制，从根源上解决了高温烧蚀问题。
\end{enumerate}

\subsection{扩孔变径方案的设计决策}

在采用一体式加工方案后，面临一个新的问题：水冷通道需要在炉体内部占据一定的径向空间（水冷孔直径0.6mm，且需保证与气流通道之间有0.5mm以上的安全壁厚）。若将全部燃料管和氧化剂管以统一孔径贯穿整个炉体，则在有限的空间内，管间距的最小值将受限于“微孔直径 + 水冷孔直径 + 安全壁厚”的综合约束。初步计算表明，若所有孔径为1mm，则管间距下限约为2.7mm。按照此参数比例，加工了一台试验炉进行点火测试，结果出现了火焰不均匀的问题，未能形成理想的平面火焰。这表明2.7mm的管间距过大，相邻射流之间的扩散混合距离过长，无法在火焰面处形成均匀的温度分布。

为兼顾“小管间距（提高火焰均匀性）”与“大水冷孔（保障高温工况下的热管理）”两个相互矛盾的设计目标，本文提出了\textbf{扩孔变径方案}：

\begin{itemize}
	\item 燃料管孔与氧化剂管孔在水冷通道下方区域（即炉体进气段至水冷层之间的区域），采用小孔径约为0.4mm，以尽可能缩小管间距，增加单位面积上的孔密度。
	\item 当微孔穿过水冷层后，在即将到达炉体顶面出口的位置，通过\textbf{30$^\circ$内倒角}将出口孔径平滑扩大。
\end{itemize}

该变径设计的流体力学合理性在于：在小孔径段，气流速度较高（出口流速约120m/s），有利于抑制回火并增强射流的穿透力；在扩口段，气流速度降低、静压恢复，有利于在火焰面处形成较为平缓的速度剖面。CFD仿真表明，该变径结构使得出口截面的速度不均匀度指数$\gamma$维持在0.97以上，未因扩口引入显著的流动畸变。

水冷通道的布局如图所示。冷却水由炉体侧面进水口进入环形集水槽，经由微孔阵列外围的轴向水冷孔（孔径4mm，共12个均匀分布于最外圈孔外侧）流至炉体顶部附近，经顶部环形汇水槽汇聚后由另一侧出水口排出。该“下进上出、环形分布”的流路设计确保了水冷层对高温区域的全覆盖。

\begin{figure}[H]
	\centering
	\begin{subfigure}{0.48\textwidth}
		\centering
		\includegraphics[width=\linewidth]{image12.png}
		\label{fig:burner_design_side}
	\end{subfigure}
	\hfill
	\begin{subfigure}{0.48\textwidth}
		\centering
		\includegraphics[width=\linewidth]{image13.png}
		\label{fig:burner_design_top}
	\end{subfigure}
	\caption{火焰炉结构侧视图与俯视图}
	\label{fig:burner_design}
\end{figure}

\subsection{三种对照方案的设计与加工}

为系统探究占孔比及绝对射流尺寸对火焰均匀性和平面度的影响，本研究设计了三种对照方案，均采用相同加工工艺制造。三种方案的核心尺寸参数如表\ref{tab:burner_designs}所示。

\begin{table}[H]
	\centering
	\caption{三种炉体设计方案对比}
	\label{tab:burner_designs}
	\begin{tabular}{ccccc}
		\toprule
		\textbf{方案} & \textbf{出口管径} & \textbf{管间距} & \textbf{间距/管径比} & \textbf{占孔比} \\
		\midrule
		炉1 & 0.8\,mm & 1.6\,mm & 2.00 & 19.63\% \\
		炉2 & 1.1\,mm & 2.2\,mm & 2.00 & 19.63\% \\
		炉3 & 1.0\,mm & 3.0\,mm & 3.00 & 8.73\% \\
		\bottomrule
	\end{tabular}
\end{table}

三种方案的对照逻辑如下：

\begin{itemize}
	\item \textbf{炉1 与 炉2（固定占孔比，考察射流绝对尺度影响）}：两方案具有完全相同的间距/管径比（2.00）和占孔比（19.63\%），主要区别在于出口管径的绝对尺寸（0.8\,mm 对比 1.1\,mm）。该对照组的目的是\textbf{在保持喷口相对分布特征不变的前提下，独立考察射流雷诺数（绝对尺寸）对燃烧混合特性的影响}。若两方案的温度场均匀性相近，则说明在此尺度范围内占孔比是主导因素，几何缩放不影响火焰性能；若存在显著差异，则说明绝对射流尺度引起的混合效果变化不可忽略，需以实测结果作为最终选型依据。
	
	\item \textbf{炉1/炉2 与 炉3（固定或近似管径，考察占孔比变化的影响）}：炉3的间距/管径比增大至3.00，对应的占孔比大幅下降至8.73\%，约为前两者的 44\%。该对照组的目的是\textbf{考察占孔比对火焰面均匀性和稳定性的影响}。占孔比越低，单位面积内的气流射流喷口越稀疏，可能导致相邻射流间的横向卷吸能力减弱，火焰面更容易出现局部空隙或温度梯度。通过对比 19.63\% 和 8.73\% 两组占孔比下的温度场扫描结果，可以判定是否存在最优化开孔率阈值。
	
	\item \textbf{综合对照意义}：通过上述两组变量控制，可以有效厘清“绝对射流尺寸”与“开孔密度（占孔比）”两个维度对火焰平面度的影响权重。三台炉子的对比实验结果将为后续标准火焰炉的定型和尺寸外推提供关键实验依据——若占孔比是决定火焰均匀性的首要因素，则应以保持高占孔比（约19.63\%）为核心设计原则；若雷诺数效应显著，则需在占孔比与射流速度之间寻求平衡点。
\end{itemize}

三种方案均采用相同的加工标准、匹配的供气面板与水冷布局，确保变量控制的严格性。加工完成后，三台炉子将依次在同一TDLAS实验台上进行相同工况（当量比\(\phi=1.0\)，质量通量 3 kg/m\(^2\)/s）的火焰温度场扫描测量，定量对比各方案的性能优劣，筛选出最优构型用于后续标定实验。
\subsection{前期点火验证}

在三种方案的正式炉体加工完成之前，课题组先后制造了两台中试炉体，开展了两个阶段的前期点火验证，以检验不同设计方案的可行性。

\textbf{第一阶段：均一孔径方案（无扩孔变径）}

第一台中试炉体按照“均一孔径”方案加工：燃料管和氧化剂管以统一孔径（1 mm）贯穿整个炉体，未采用扩孔变径结构。受限于水冷通道的布置需求，该方案的管间距被设置为2.7 mm。点火测试采用H$_2$/O$_2$燃烧工况，氢气流量150 SLM、氧气流量75 SLM（质量通量约1.25 kg/m$^2$/s，当量比$\phi = 1.0$）。

测试结果如下：

\begin{itemize}
	\item 火焰面呈现明显的\textbf{不均匀结构}：炉体中心区域火焰亮度较高，边缘区域亮度偏低，未能形成理想的平面火焰形态。
	\item 局部出现\textbf{火焰抬升}和\textbf{分离}现象，部分孔位的射流未能在出口处及时着火，导致火焰面断裂。
\end{itemize}

该阶段验证表明：当管间距增大至2.7 mm时，相邻射流之间的扩散混合距离过大，剪切层在火焰面高度处尚未充分发展，无法形成连续均匀的反应面。这一结果为扩孔变径方案的提出提供了直接的实验依据——必须在保障水冷能力的前提下尽可能缩小管间距。

\textbf{第二阶段：扩孔变径方案（方案1参数）}

第二台中试炉体按照方案1的参数加工：采用扩孔变径结构（0.4 mm/0.8 mm变径，30$^\circ$内倒角），管间距缩小至1.6 mm。水冷通道保持与第一阶段相同的布局和流量（5 L/min，进水温度25$^\circ$C）。点火工况与第一阶段完全一致（H$_2$ 150 SLM，O$_2$ 75 SLM，$\phi = 1.0$），连续运行1小时后：

\begin{itemize}
	\item 水冷系统工作正常，出水口温度较进水口温升约8$^\circ$C，未出现局部过热或冷却水沸腾现象。
	\item 火焰面呈均匀的\textbf{淡蓝色平面火焰}，肉眼观察无局部高温区或火焰抬升现象，火焰面平整连续，覆盖整个炉体出口区域。
	\item 运行结束后拆解检查，炉体顶部未见烧蚀痕迹或变色，微孔出口边缘完整无损，表明水冷结构有效地将炉体温度控制在材料安全阈值以下。
\end{itemize}

上述两阶段验证形成了鲜明的对比：均一孔径大间距方案（2.7 mm）未能产生均匀的平面火焰，而扩孔变径小间距方案（1.6 mm）实现了均匀稳定的平面火焰。该对比清晰地表明：\textbf{缩小的管间距是获得均匀平面火焰的关键，而扩孔变径设计是同时实现小间距与水冷集成的有效技术路径}。

\begin{figure}[H]
	\centering
	\begin{subfigure}{0.45\textwidth}
		\centering
		\includegraphics[width=\linewidth]{image18.png}
		\label{fig:flame_uneven}
	\end{subfigure}
	\hfill
	\begin{subfigure}{0.45\textwidth}
		\centering
		\includegraphics[width=\linewidth]{image17.png}
		\label{fig:flame_even}
	\end{subfigure}
	\caption{不同方案火焰形态对比}
	\label{fig:flame_validation}
\end{figure}


\section{数值仿真方法}

本研究采用数值仿真与实验验证相结合的方法，从热力学平衡和冷态流体动力学两个层面对不同炉体设计方案进行评估。热力学仿真用于确定不同燃料在不同质量通量下的绝热火焰温度和热流密度，为实验工况设计提供理论依据；冷态流体仿真用于评估炉体出口截面的气流混合均匀性，为加工方案筛选提供依据。

\subsection{热力学平衡仿真（Cantera）}

热力学平衡仿真采用\textbf{Cantera}开源化学反应动力学软件包（版本2.6.0），调用其\texttt{IdealGasConstPressureReactor}定压反应器模型，分别对H$_2$/O$_2$和C$_2$H$_2$/O$_2$两种燃料组合进行化学平衡计算。

计算采用\textbf{GRI-Mech 3.0}化学动力学机理（包含53种组分和325个基元反应），该机理在碳氢燃料及氢气燃烧模拟中具有广泛的适用性和验证基础。对于H$_2$/O$_2$体系，GRI-Mech 3.0中的H$_2$/O$_2$子机理已针对点火延迟和火焰传播速度进行了系统校验，可较好地预测高温条件下的平衡组分和绝热火焰温度。

计算条件设置如下：

\begin{itemize}
	\item 初始温度：300 K
	\item 初始压力：1 atm（绝热等压燃烧，最终压力保持1 atm）
	\item 当量比：$\phi = 1.0$（化学恰当比）
	\item 质量通量范围：0.5--3.0 kg/m$^2$/s（根据炉体设计工况确定）
	\item 燃料/氧化剂配比：依据化学恰当比的H$_2$/O$_2$（2:1摩尔比）和C$_2$H$_2$/O$_2$（2:5摩尔比）分别设定
\end{itemize}

计算输出参数包括：绝热火焰温度（Adiabatic Flame Temperature）、平衡组分摩尔分数（H$_2$O、CO$_2$、CO、OH、H、O等）、焓差（反应物与产物之间的总焓变化）以及表面热流密度（根据质量通量与焓差的乘积换算）。

\subsubsection{H$_2$/O$_2$火焰计算结果}

H$_2$/O$_2$火焰（当量比$\phi=1.0$）在不同质量通量下的热力学参数如表\ref{tab:cantera_h2}所示。

\begin{table}[H]
	\centering
	\caption{Cantera计算：H$_2$/O$_2$火焰热力学参数（$\phi=1.0$，1 atm）}
	\label{tab:cantera_h2}
	\begin{tabular}{ccccc}
		\toprule
		质量通量 (kg/m$^2$/s) & 焓差 (J/m$^2$/s) & 热流密度 (MW/m$^2$) & 温度 (K) \\
		\midrule
		0.5 & 3,771,483 & 6.5 & 2456 \\
		1 & 7,542,967 & 12.0 & 2643 \\
		1.5 & 11,314,450 & 17.5 & 2754 \\
		2 & 15,085,934 & 22.0 & 2835 \\
		3 & 22,628,901 & 30.0 & 2953 \\
		\bottomrule
	\end{tabular}
\end{table}

由表\ref{tab:cantera_h2}可知，H$_2$/O$_2$火焰的绝热温度随质量通量增加而升高，从0.5 kg/m$^2$/s时的2456 K上升至3.0 kg/m$^2$/s时的2953 K，但升温速率逐渐减缓（接近H$_2$/O$_2$的理论绝热温度上限$\sim$3070 K）。热流密度从6.5 MW/m$^2$增至30.0 MW/m$^2$，呈近似线性增长趋势，表明高质量通量工况下炉体顶部承受的热负荷急剧增加，对水冷系统的散热能力提出了更高要求。

\subsubsection{C$_2$H$_2$/O$_2$火焰计算结果}

C$_2$H$_2$/O$_2$火焰（当量比$\phi=1.0$）在不同质量通量下的热力学参数如表\ref{tab:cantera_c2h2}所示。

\begin{table}[H]
	\centering
	\caption{Cantera计算：C$_2$H$_2$/O$_2$火焰热力学参数（$\phi=1.0$，1 atm）}
	\label{tab:cantera_c2h2}
	\begin{tabular}{ccccc}
		\toprule
		质量通量 (kg/m$^2$/s) & 焓差 (J/m$^2$/s) & 热流密度 (MW/m$^2$) & 温度 (K) \\
		\midrule
		0.5 & 2,308,730 & 5.35 & 2591 \\
		1 & 4,617,460 & 9.80 & 2768 \\
		1.5 & 6,926,190 & 13.70 & 2874 \\
		2 & 9,234,920 & 17.14 & 2951 \\
		3 & 13,852,380 & 22.93 & 3064 \\
		\bottomrule
	\end{tabular}
\end{table}

C$_2$H$_2$/O$_2$火焰的绝热温度在相同质量通量下普遍高于H$_2$/O$_2$火焰，在3.0 kg/m$^2$/s时达到3064 K，已接近乙炔/氧气理论绝热温度上限（$\sim$3100 K）。表明C$_2$H$_2$/O$_2$体系更适用于需要更高标定温度的极端工况场景。

H$_2$/O$_2$火焰的Flame Spectrum平台燃烧仿真结果如图所示。

\begin{figure}[H]
	\centering
	\begin{subfigure}{0.48\textwidth}
		\centering
		\includegraphics[width=\linewidth]{image19.png}
		\label{fig:c2h2_temp}
	\end{subfigure}
	\hfill
	\begin{subfigure}{0.48\textwidth}
		\centering
		\includegraphics[width=\linewidth]{image20.png}
		\label{fig:c2h2_species}
	\end{subfigure}
	\caption{H$_2$/O$_2$火焰燃烧仿真结果}
	\label{fig:c2h2_sim}
\end{figure}

\subsubsection{仿真结果讨论}

对比两组计算结果可得出以下结论：

第一，两种燃料体系的热流密度均随质量通量线性增加，其中H$_2$/O$_2$体系在高质量通量下的热流密度（30 MW/m$^2$）显著高于C$_2$H$_2$/O$_2$体系（22.93 MW/m$^2$），尽管前者的绝热温度略低。该差异源于H$_2$/O$_2$反应释放的化学热中更大的比例转化为显热而非离解能，导致更多的热量传递至炉体表面。

第二，在质量通量2--3 kg/m$^2$/s范围内，两种燃料的绝热温度均超过2800 K，远超传统NPL标准火焰（丙烷/空气体系）的2260 K上限。该结果表明，采用H$_2$/O$_2$或C$_2$H$_2$/O$_2$体系配合水冷炉体设计，具备将火焰标定温度范围拓展至3000 K的理论潜力。

第三，上述热力学计算结果为后续实验工况的选取提供了关键依据：低质量通量（0.5--1 kg/m$^2$/s）适用于验证性实验和系统调试，中等质量通量（1.5--2 kg/m$^2$/s）为常规标定工况，高质量通量（3 kg/m$^2$/s及以上）为目标高温工况。

\subsection{冷态流体仿真验证}

前述热力学分析确定了炉体所需承受的理论热负荷和温度水平，但炉体能否在实验条件下产生均匀的平面火焰，取决于其出口截面的气流混合均匀性。为此，采用Star-CCM+开展冷态流体仿真验证。

本研究采用Siemens Industry旗下CD-adapco公司开发的\textbf{Star-CCM+ 2024.10}作为主要计算流体动力学（CFD）仿真平台。仿真工作覆盖几何建模、网格划分、物理模型选取、边界条件设定、求解及后处理全流程，旨在定量评估不同炉体设计方案在冷态条件下的气流混合均匀性，为最终加工方案提供依据。

\subsubsection{几何模型与计算域设定}

以实验用非预混平面绝热火焰炉的实际尺寸为基准，建立三维几何模型。为定量评估前述三种设计方案（火焰炉1、火焰炉2、火焰炉3）在冷态条件下的流动均匀性差异，分别建立与各方案几何参数精确对应的三维计算模型。三种方案的核心几何参数如表\ref{tab:burner_designs}所示，均采用扩孔变径结构，微孔段长度$L=15$mm。

\begin{itemize}
	\item \textbf{火焰炉1}：管间距$s=1.6$mm，氢气管径$0.4/0.8$mm（变径），氧气管径$0.4/0.8$mm（变径），扩张段倒角$30^\circ$，水冷结构有。
	\item \textbf{火焰炉2}：管间距$s=2.2$mm，氢气管径$0.4/1.1$mm（变径），氧气管径$0.4/1.1$mm（变径），扩张段倒角$30^\circ$，水冷结构有。
	\item \textbf{火焰炉3}：管间距$s=3.0$mm，氢气管径$0.4/0.8$mm（变径），氧气管径$0.4/0.8$mm（变径），扩张段倒角$30^\circ$，水冷结构有。
\end{itemize}

为进一步捕捉火焰面附近的流动发展，在微孔出口上方额外延伸\textbf{10 mm}的自由射流区域（环境大气域），该区域直径取50 mm以消除边界约束对中心射流的影响。计算域底部（微孔入口截面）设定为质量流量入口边界，顶部及侧面设定为压力出口边界，在出口上方4 mm处查看气流混合的均匀性。

\subsubsection{网格划分策略}

网格生成采用Star-CCM+内置的\textbf{多面体网格（Polyhedral Mesh）}与\textbf{棱柱层网格（Prism Layer Mesh）}相结合的策略，具体参数如下：

\begin{itemize}
	\item \textbf{体网格类型}：多面体网格。相较于四面体网格，多面体网格在相同计算精度条件下所需的网格数量更少（约减少40\%--60\%），且对复杂几何的适应性更强，收敛性更好。
	\item \textbf{近壁面处理}：在微孔壁面及出口扩张段等固体壁面附近，生成\textbf{5层棱柱层网格}，厚度增长率设为1.2，以确保壁面边界层内的速度梯度被充分解析。
	\item \textbf{边界层分辨率}：控制第一层网格高度使得壁面$y^+ < 1$，满足低雷诺数湍流模型对近壁面网格精度的要求。
	\item \textbf{局部加密}：在\textbf{微孔出口区域（直径2 mm范围内）}以及\textbf{射流剪切层（沿射流方向延伸10 mm、径向向外扩展1 mm的带状区域）}进行额外网格加密。加密策略采用Star-CCM+的\textbf{体网格控制（Mesh Refinement Region）}功能，设定目标网格尺寸为0.02 mm（约为微孔直径的2\%--5\%），以确保剪切层内大梯度物理量（速度、湍动能）被准确捕捉。
\end{itemize}

\subsubsection{物理模型与湍流模型选择}

基于仿真目标为冷态等温流动（无化学反应、无热源），激活以下物理模型：

\begin{itemize}
	\item \textbf{空间维度}：三维（3D）
	\item \textbf{时间模式}：稳态（Steady State）
	\item \textbf{材料属性}：气体（理想气体，密度由理想气体状态方程计算，以考虑可压缩性效应）
	\item \textbf{流动模型}：层流（基于微孔出口流速和特征尺寸计算得到的雷诺数$Re \approx 500$--1000，处于层流或转捩区，层流假设可合理反映低雷诺数射流的混合特性）
\end{itemize}

\subsubsection{边界条件设置}

计算域各边界条件设定如下表\ref{tab:simulation_bc}所示：

\begin{table}[H]
	\centering
	\caption{Star-CCM+仿真边界条件汇总}
	\label{tab:simulation_bc}
	\begin{tabular}{p{4cm}p{4.5cm}p{6cm}}
		\toprule
		\textbf{边界类型} & \textbf{位置} & \textbf{具体设置} \\
		\midrule
		质量流量入口（Mass Flow Inlet） & 微孔底部（H$_2$孔，O$_2$孔） & 依据设计工况设定：质量通量4 kg/m$^2$/s，当量比$\phi=1.0$，氧气流量50 SLM，氢气流量100 SLM。 \\
		压力出口（Pressure Outlet） & 计算域顶部及侧面 & 静压设为0 Pa（表压），温度300 K。 \\
		对称面（Symmetry Plane） & 周期性单元两侧 & 法向速度为零，所有物理量法向梯度为零。 \\
		无滑移壁面（No-Slip Wall） & 微孔内壁面及炉体顶部表面 & 速度为零，绝热壁面条件（冷态仿真）。 \\
		\bottomrule
	\end{tabular}
\end{table}

\subsubsection{仿真结果}

在设定工况下（质量通量4 kg/m$^2$/s，当量比$\phi=1.0$），三种方案的冷态仿真结果对比如表\ref{tab:simulation_comparison}所示。

\begin{table}[H]
	\centering
	\caption{三种方案冷态仿真结果对比}
	\label{tab:simulation_comparison}
	\begin{tabular}{cccccc}
		\toprule
		\textbf{方案} & \textbf{出口管径 (mm)} & \textbf{孔间距 (mm)} & \textbf{间距/管径比} & \textbf{速度均匀性指数 $\gamma$} & \textbf{均匀性判断} \\
		\midrule
		火焰炉1 & 0.8 & 1.6 & 2.00 & 0.971 & 均匀 \\
		火焰炉2 & 1.1 & 2.2 & 2.00 & 0.964 & 均匀 \\
		火焰炉3 & 0.8 & 3.0 & 3.75 & 0.833 & 不均匀 \\
		\bottomrule
	\end{tabular}
\end{table}

仿真结果表明：

\begin{enumerate}
	\item \textbf{火焰炉1与火焰炉2均满足均匀性要求}：两者的速度均匀性指数$\gamma$分别为0.971和0.964，均超过设计目标值0.95。火焰炉1的均匀性略优于火焰炉2（0.971 vs 0.964），表明在保持间距/管径比恒定的条件下，整体尺度放大对流动均匀性无负面影响。
	
	\item \textbf{火焰炉3均匀性显著下降}：当管间距扩大至3.0 mm而出口管径保持0.8 mm不变时，间距/管径比由2.00增大至3.75，速度均匀性指数$\gamma$降至0.833，低于0.95的设计目标。该结果表明，当间距/管径比超过约2.0--2.5的临界范围后，相邻射流之间的间距过大，剪切层在火焰面高度处无法充分发展并相互融合，导致出口截面出现明显的速度不均匀分布。
	
	\item \textbf{间距/管径比是决定均匀性的关键无量纲参数}：火焰炉1与火焰炉2虽然绝对尺寸不同（出口管径0.8 mm vs 1.1 mm，间距1.6 mm vs 2.2 mm），但二者间距/管径比均为2.00，仿真结果显示二者的流动均匀性处于同一水平（$\gamma$分别为0.971和0.964）。相比之下，火焰炉3的间距/管径比增大至3.75，均匀性即显著恶化。这一发现证实，在本文所研究的尺度范围内，影响出口截面速度均匀性的主导因素是\textbf{间距与管径的比值}，而非单一几何参数的绝对大小。
\end{enumerate}

\begin{figure}[H]
	\centering
	\begin{subfigure}{0.32\textwidth}
		\centering
		\includegraphics[width=\linewidth]{image14.png}
		\label{fig:sim_burner1}
	\end{subfigure}
	\hfill
	\begin{subfigure}{0.32\textwidth}
		\centering
		\includegraphics[width=\linewidth]{image15.png}
		\label{fig:sim_burner2}
	\end{subfigure}
	\hfill
	\begin{subfigure}{0.32\textwidth}
		\centering
		\includegraphics[width=\linewidth]{image16.png}
		\label{fig:sim_burner3}
	\end{subfigure}
	\caption{三种方案出口截面速度均匀性仿真对比}
	\label{fig:sim_comparison}
\end{figure}

上述仿真结论与前期点火验证的实验结果形成了有效印证：第一台中试炉体（均一孔径，管间距2.7 mm，管径1.0 mm，间距/管径比$2.7$）火焰不均匀；第二台中试炉体（扩孔变径方案1，管间距1.6 mm，出口管径0.8 mm，间距/管径比$2.00$）火焰均匀。仿真与实验共同表明：当间距/管径比控制在约2.0左右时，可保证良好的流动均匀性与火焰平面度。

综合热力学仿真与冷态流体仿真结果，确定火焰炉1（$s=1.6$mm，$d_{exit}=0.8$mm，间距/管径比$=2.00$）作为后续标定实验的标准构型。该方案在保障火焰均匀性的同时，具有最小的管间距和最高的孔密度，可为TDLAS测量提供最优的空间分辨率；火焰炉2的仿真结果也验证了在保持间距/管径比恒定的条件下几何缩放的可行性，为未来炉体尺寸的系列化设计提供了理论依据。

\subsection{供气与流量控制}

本研究采用\textbf{临界流小孔（限流孔）}作为流量控制的核心元件。其基本原理为：当气体流经小孔截面时，若小孔上下游压比低于临界压比，气流在小孔喉部达到声速（Ma=1），此时流经小孔的质量流量仅取决于上游滞止压力与温度，与下游压力波动完全解耦。

临界压比$\Pi_{cr}$由等熵流动关系确定：
\begin{equation}
	\Pi_{cr} = \frac{p_{down}}{p_{up}} = \left( \frac{2}{\gamma + 1} \right)^{\frac{\gamma}{\gamma - 1}}
	\label{eq:critical_pressure}
\end{equation}
其中$\gamma$为气体的比热比。对于双原子气体（H$_2$、O$_2$、N$_2$），$\gamma \approx 1.4$，计算得$\Pi_{cr} \approx 0.528$；对于多原子气体（CH$_4$），$\gamma \approx 1.32$，$\Pi_{cr} \approx 0.542$。只要小孔上下游压比小于上述临界值，即可保证限流孔工作在声速状态。

在临界状态下，通过小孔的质量流量$\dot{m}$由以下公式给出：
\begin{equation}
	\dot{m} = C_d \cdot A \cdot \frac{p_{up}}{\sqrt{T_{up}}} \cdot \sqrt{\frac{\gamma}{R} \left( \frac{2}{\gamma + 1} \right)^{\frac{\gamma + 1}{\gamma - 1}}}
	\label{eq:choked_flow}
\end{equation}
其中：$C_d$为流量系数（对于锐边薄壁小孔，$C_d \approx 0.72$）；$A = \pi d^2 / 4$为小孔截面积；$p_{up}$和$T_{up}$分别为小孔上游的滞止压力和温度；$R$为气体常数。

根据公式\eqref{eq:choked_flow}，在给定目标流量$\dot{m}$和上游滞止压力$p_{up}$的条件下，所需小孔直径$d$可反算得到：
\begin{equation}
	d = \sqrt{\frac{4\dot{m}}{C_d \pi p_{up}} \cdot \sqrt{\frac{R T_{up}}{\gamma}} \cdot \left( \frac{\gamma + 1}{2} \right)^{\frac{\gamma + 1}{2(\gamma - 1)}}}
	\label{eq:orifice_diameter}
\end{equation}

在当量比$\phi=1$、质量通量4 kg/m$^2$/s、上游压力0.3 MPa、室温$T_{up}=300$ K条件下，按上述公式计算可得H$_2$通道限流孔径约2.04 mm，O$_2$通道限流孔径约2.88 mm。不同质量通量下的流量及对应限流孔径如表\ref{tab:flow_orifice}所示。

小孔限流方案相较于直接使用MFC具有以下优势：第一，将流量控制从“MFC直接调节流量”转变为“MFC调节上游压力+固定小孔限流”，MFC始终工作在中高量程区间，避免了小流量工况下的精度下降；第二，火焰燃烧时炉体压力波动不影响通过小孔的流量，提升了燃烧稳定性；第三，小孔为纯机械元件，长期重复性优于电子流量计，且便于更换以快速切换工况。

\begin{table}[H]
	\centering
	\caption{不同质量通量下的流量与限流孔径（H$_2$/O$_2$，$\phi=1$，上游压力0.3 MPa）}
	\label{tab:flow_orifice}
	\begin{tabular}{cccccc}
		\toprule
		质量通量 & 质量流量 & H$_2$流量 & O$_2$流量 & H$_2$孔径 & O$_2$孔径 \\
		(kg/m$^2$/s) & (kg/s) & (SLM) & (SLM) & (mm) & (mm) \\
		\midrule
		0.5 & $8.0\times10^{-4}$ & 60 & 30 & 0.72 & 1.02 \\
		1 & $1.6\times10^{-3}$ & 120 & 60 & 1.02 & 1.44 \\
		1.5 & $2.4\times10^{-3}$ & 180 & 90 & 1.25 & 1.77 \\
		2 & $3.2\times10^{-3}$ & 240 & 120 & 1.44 & 2.04 \\
		3 & $4.8\times10^{-3}$ & 360 & 180 & 1.77 & 2.50 \\
		4 & $6.4\times10^{-3}$ & 480 & 240 & 2.04 & 2.88 \\
		\bottomrule
	\end{tabular}
\end{table}

供气面板所有管路均采用1/2英寸EP级不锈钢管，接头为SMC双卡套密封形式，关键节点配置防回火逆止阀，确保氢气/氧气供气的本质安全性。


\section{伺服运动平台}

二维高精度扫描运动平台是TDLAS二维温度场测量的关键执行机构，负责将火焰炉相对于激光束进行精确的二维步进位移，实现逐点空间分辨测量。

\subsection{硬件选型与系统架构}

运动机构选用THK KR-20型高精度线性模组，采用滚珠丝杠传动配合直线导轨导向，重复定位精度$\pm 0.01$ mm，行程250 mm $\times$ 250 mm，最大负载能力20 kg，足以承载火焰炉及配套的冷却水管路和气路。双轴采用正交堆叠布局，Y轴模组（垂直方向，对应火焰高度）固定于光学平台，X轴模组（水平方向，对应火焰径向）安装于Y轴滑块之上，火焰炉固定在X轴滑块上，实现炉体相对于固定激光束的二维移动。

驱动与控制采用\textbf{CC-D3P040A2-A型伺服驱动器}配套\textbf{CC-M3H040-NN14型伺服电机}。驱动器采用RS-485串行通讯接口与上位机连接，支持Modbus RTU协议，可通过LabVIEW编写控制程序实现位置、速度、加速度的全参数可编程控制。控制架构如图所示。

\begin{figure}[H]
	\centering
	\begin{subfigure}{0.45\textwidth}
		\centering
		\includegraphics[width=\linewidth]{image41.png}
		\label{fig:servo_control_front}
	\end{subfigure}
	\hfill
	\begin{subfigure}{0.45\textwidth}
		\centering
		\includegraphics[width=\linewidth]{image32.png}
		\label{fig:servo_control_program}
	\end{subfigure}
	\caption{LabVIEW伺服平台前面板与控制程序界面}
	\label{fig:servo_control}
\end{figure}

\subsection{通讯协议与控制参数配置}

在通过LabVIEW实现对伺服电机的控制前，需完成以下关键配置：

\textbf{1. 控制器地址设置}
当使用多个控制器（例如需要同时控制X轴和Y轴两台电机）时，必须为每个控制器分配唯一的通讯地址。控制器上的地址拨码开关设定地址位后，LabVIEW上位机发送的控制指令中需包含对应的地址字节，驱动程序根据地址字节区分指令的接收对象。地址位发生变化时，指令帧最后的CRC校验码需重新计算。

\textbf{2. 使能控制与通讯模式切换}
伺服驱动器通常支持I/O控制和RS-485通讯控制两种模式。在单独使用LabVIEW进行控制而不开启上位机调试软件时，除了发送使能开指令（Enable）外，还需额外写入控制模式切换指令，将驱动器的控制源从I/O端子切换至RS-485通讯端口。否则驱动器虽已使能，但对LabVIEW发送的位移指令无响应。

\textbf{3. 电子齿轮比设置}
电机的实际移动量与驱动器接收到的脉冲数之间的关系由电子齿轮比决定。以本系统使用的电机为例，默认齿轮比下，发送1000个脉冲仅对应0.1圈旋转。若需要在单次指令下旋转完整一圈，则需发送10000个脉冲或重新设置齿轮比参数。通过修改齿轮比分母（由默认的2500改为1250），可使1000个脉冲对应的旋转圈数从1/10圈提升至1/5圈。电子齿轮比的合理设置直接影响扫描步长的精度和系统响应速度，需根据实际扫描步长（0.5 mm）和丝杠导程折算后确定最优比值。

\textbf{4. 转速与加减速设置}
设置目标运动速度（单位：r/min）和加速度/减速度参数，使电机在启动和停止时平滑过渡，避免因急停引入的机械振动影响光路对准。合理设置加减速时间还有利于缩短扫描周期、提高实验效率。

\subsection{LabVIEW控制程序设计}

本项目的伺服控制程序基于LabVIEW开发，核心功能模块包括：

\begin{enumerate}
	\item \textbf{VISA串口通讯模块}：通过RS-485总线与驱动器通讯，配置波特率（默认19200 bps）、数据位（8位）、停止位（1位）、校验位（无），实现指令下发与状态读取。
	
	\item \textbf{位置指令生成模块}：根据预设的扫描矩阵（X方向$-15$至$+15$ mm，Y方向1至5 mm），自动生成有序的绝对位置指令序列，并实时转换为电机所需的脉冲数或对应通讯指令。
	
	\item \textbf{位置状态读取与闭环确认模块}：在每次移动指令执行完成后，读取驱动器的当前位置反馈值（从反馈帧的第4位开始截取4位数据，强制转换为10进制整数），与目标位置进行比对，确认到位后方可向数据采集系统发送“采集触发”信号。该流程保证了光谱数据与空间位置的一一对应，避免因定位误差引入的温度场重建畸变。
	
	\item \textbf{多电机轮询控制}：当需要同时控制X轴和Y轴两台电机时，将两台驱动器的RS-485通讯线（A/B口）并联接入同一COM口，通过不同的地址位区分控制对象。LabVIEW程序采用轮询方式——在一个循环中交替发送X轴和Y轴的问询指令，驱动器根据地址位选择性响应，实现单串口总线上的多轴分时控制。轮询周期由while循环中的迭代计数与2取余判断奇偶来确定当前时刻访问的轴号，循环内两个分支分别对应X轴和Y轴的指令下发与状态读取。
\end{enumerate}

\subsection{扫描策略与数据采集联动}

扫描策略采用\textbf{逐行光栅扫描模式}：Y方向（高度）从火焰面底部1 mm开始，步长0.5 mm，逐层升高至5 mm，共9层；X方向（水平）在每个Y高度层上从$-15$ mm至$+15$ mm，步长0.5 mm，共61个点。整个二维场总计$61 \times 9 = 549$个测量点，每个测量点包含100次波长扫描的累加平均。

具体联动流程如下：

\begin{enumerate}
	\item 上位机LabVIEW程序通过RS-485指令将伺服平台移动至目标位置（$X_i$, $Y_j$）；
	\item 到位后程序自动触发NI USB-6361数据采集卡开始采集——同步记录测量光路、标准具干涉光路及背景光路的电压信号，采样率50 MS/s，单次采集时长覆盖完整的波长扫描周期；
	\item 同一位置完成300次扫描平均后，LabVIEW程序控制平台移动至下一位置（$X_{i+1}$, $Y_j$），重复上述步骤；
	\item 完成一整行（X方向全行程）后，Y方向步进0.5 mm至下一高度层，X方向反向继续扫描（蛇形路径以减少空行程时间）。
\end{enumerate}

整场测量共需采集$549 \times 100 = 54{,}900$条光谱，总耗时约1.5小时（含平台移动和通讯稳定时间）。采集过程中，所有原始数据（包含测量光路、标准具光路、背景光路的三路信号）与对应空间坐标（$X_i$, $Y_j$）一起存储为结构化数据文件，供后续数据处理流水线调用。